\subsection{Explicit Runge--Kutta Implementation}
\label{subsec:rk_implementation}

The explicit Runge--Kutta methods are implemented using a
common solver function, \texttt{\_solve\_with\_explicit\_rk}.
The method is specified through the coefficients
\texttt{rk\_a} and \texttt{rk\_b}, which determine the
intermediate stages and the final weighted update,
respectively. This allows the same implementation to be used
for both the second-order explicit midpoint method and the
classical fourth-order Runge--Kutta method derived in the
previous section.

\subsubsection{Stage Evaluation and Temperature Update}

At each pseudo-time iteration, the solver evaluates the
right-hand-side operator defined in
Equation~\eqref{eq:rk_rhs_operator} at several intermediate
temperature fields. For a method with $m$ stages, the
intermediate fields and stage evaluations are expressed
generally as
\begin{align}
T^{(s)}
&=
T^n+\Delta t
\sum_{r=1}^{s-1}a_{sr}k_r,
\qquad s=1,\ldots,m,
\label{eq:rk_implementation_stages}
\\
k_s
&=
\mathcal{F}\left(T^{(s)}\right),
\qquad s=1,\ldots,m,
\label{eq:rk_implementation_rhs}
\\
T^{n+1}
&=
T^n+\Delta t
\sum_{s=1}^{m}b_s k_s,
\label{eq:rk_implementation_update}
\end{align}
where $a_{sr}$ and $b_s$ are the Runge--Kutta coefficients,
and $T^{(1)}=T^n$. These general expressions correspond to
the specific stage equations and updates given in
Equations~\eqref{eq:rk2_k1}--\eqref{eq:rk2_update} for RK2
and Equations~\eqref{eq:rk4_k1}--\eqref{eq:rk4_update}
for RK4.

The implementation advances only the physical-cell
temperatures explicitly. At each stage, a copy of the
current temperature field is created, and the intermediate
stage contributions are added to its physical cells.
The temperature boundary conditions are then reapplied
before evaluating the spatial right-hand side. This ensures
that the ghost-cell values are consistent with the boundary
conditions at every stage.

After all stages have been evaluated, their right-hand sides
are combined using the coefficients $b_s$ to update the
physical-cell temperatures. The boundary conditions are
reapplied once more to obtain a consistent updated
temperature field.

\begin{algorithm}[H]
\caption{One Explicit Runge--Kutta Time Step}
\label{alg:rk_single_step}
\begin{algorithmic}[1]
\Require Temperature field $T^n$, time step $\Delta t$,
Runge--Kutta coefficients $a_{sr},b_s$, velocity fields
$u,v$, and spatial-operator parameters
\Ensure Updated temperature field $T^{n+1}$

\State Apply the temperature boundary conditions to $T^n$
\State Initialize an empty list of stage right-hand sides

\For{$s=1,\ldots,m$}
    \State Set $T^{(s)}\gets T^n$
    \If{$s>1$}
        \State Construct the intermediate stage using
        Equation~\eqref{eq:rk_implementation_stages}
        using only physical-cell values
    \EndIf
    \State Apply the temperature boundary conditions to $T^{(s)}$
    \State Evaluate $k_s$ using
    Equation~\eqref{eq:rk_implementation_rhs}
    \State Store $k_s$ in the stage right-hand-side list
\EndFor

\State Initialize the total increment to zero
\For{$s=1,\ldots,m$}
    \State Add $b_s k_s$ to the total increment
    over physical cells
\EndFor
\State Update the physical-cell temperatures using
Equation~\eqref{eq:rk_implementation_update}
\State Apply the temperature boundary conditions to $T^{n+1}$
\State \Return $T^{n+1}$
\end{algorithmic}
\end{algorithm}

For the explicit midpoint method, the solver uses two stages,
with coefficients $a_{21}=1/2$ and $b_1=0$, $b_2=1$.
For classical RK4, four stages are used, with
\[
\begin{gathered}
a_{21}=\frac12,\qquad
a_{32}=\frac12,\qquad
a_{43}=1,\\
b_1=\frac16,\qquad
b_2=\frac13,\qquad
b_3=\frac13,\qquad
b_4=\frac16.
\end{gathered}
\]
Thus, the two methods differ only in their stage coefficients
and the number of right-hand-side evaluations required per
iteration.

\subsubsection{Pseudo-Time Marching and Convergence}

The solver repeatedly advances the temperature field by one
Runge--Kutta step until the steady-state residual reaches the
specified tolerance or the maximum number of iterations is
exceeded. Before iteration begins, the initial temperature
field is checked for compatibility with the expected array
dimensions and the temperature boundary conditions are
applied.

After each complete time step, the RMS residual is evaluated
using Equation~\eqref{eq:steady_residual_rms}. The iteration
number and residual are recorded, and temperature snapshots
are stored at the specified interval. The initial field is
stored at iteration zero, and the final field is included
even when the final iteration does not coincide with a
snapshot interval.

The convergence criterion is
\begin{equation}
\boxed{
R_{\mathrm{RMS}}\leq\varepsilon,
\qquad \varepsilon=10^{-12}.
}
\label{eq:rk_convergence_criterion}
\end{equation}
The computation time is measured over the pseudo-time
iteration loop. If the residual tolerance is not reached
within the prescribed maximum number of iterations, a
warning is issued and the final computed solution and
diagnostic histories are returned.

\begin{algorithm}[H]
\caption{Runge--Kutta Pseudo-Time Marching to Steady State}
\label{alg:rk_complete_solver}
\begin{algorithmic}[1]
\Require Initial temperature $T^0$, time step $\Delta t$,
Runge--Kutta coefficients $a_{sr},b_s$, residual tolerance
$\varepsilon$, maximum iterations $N_{\max}$, and snapshot
interval $N_s$
\Ensure Final temperature field $T$, iteration and residual
histories, computation time, and temperature snapshots

\State Validate the time step, iteration limit, snapshot interval,
stage coefficients, and initial field dimensions
\State Apply the temperature boundary conditions to $T^0$
\State Initialize iteration and residual histories as empty
\State Store $T^0$ as the initial snapshot at iteration $0$
\State Start the computation timer
\State $\mathrm{converged}\gets\mathrm{false}$

\For{$n=1,\ldots,N_{\max}$}
    \State Advance $T$ by one Runge--Kutta step using
    Algorithm~\ref{alg:rk_single_step}
    \State Evaluate $R_{\mathrm{RMS}}$ using
    Equation~\eqref{eq:steady_residual_rms}
    \State Append $n$ to the iteration history
    \State Append $R_{\mathrm{RMS}}$ to the residual history

    \If{$n\bmod N_s=0$}
        \State Store the current temperature field as a snapshot
        \State Record snapshot iteration $n$
    \EndIf

    \If{$R_{\mathrm{RMS}}\leq\varepsilon$}
        \State $\mathrm{converged}\gets\mathrm{true}$
        \State \textbf{break}
    \EndIf
\EndFor

\State Stop the timer and calculate computation time

\If{the final iteration is not already stored as a snapshot}
    \State Store the final temperature field and iteration
\EndIf

\If{not converged}
    \State Issue a warning that the residual tolerance was
    not reached within $N_{\max}$ iterations
\EndIf

\State \Return $T$, iteration history, residual history,
computation time, temperature snapshots, and snapshot iterations
\end{algorithmic}
\end{algorithm}

The RK2 and RK4 implementations therefore share the same
boundary-condition treatment, residual evaluation, convergence
criterion, and snapshot procedure. Their main differences
are the number of stages and the associated Runge--Kutta
coefficients. Both methods are explicit and conditionally
stable; consequently, the pseudo-time step must be selected
to maintain numerical stability while allowing efficient
convergence to the steady solution.